\documentclass[10pt,a4paper]{amsart}
\usepackage[margin=19mm]{geometry}
\usepackage{amsmath,amssymb,mathtools}
\usepackage[scr=rsfs,cal=euler]{mathalfa}
\usepackage{xfrac}
\usepackage[unicode,hidelinks,bookmarks=false]{hyperref}
\newcommand{\ie}{i.e.\ }
\newcommand{\cf}{cf.\ }
\newcommand{\resp}{respectively\ }
\newcommand{\Subsection}{\subsection}
\providecommand{\nobreakdash}{\nobreak\hbox{-}}
\setlength{\emergencystretch}{2em}
\multlinegap0pt
\usepackage{amssymb}
\usepackage{mathtools}
\usepackage[shortlabels]{enumitem}
\usepackage{xcolor}
\usepackage{float}
\usepackage{algorithm}\usepackage{algpseudocode}
\newtheorem{proposition}{Proposition}
\title{Multi-Objective Planning for Healthcare Facility Resilience: Mitigate Now or Respond Later?}
\author{Gizem Toplu-Tutay, John J. Hasenbein, Matt Kammer-Kerwick, Erhan Kutanoglu}
\date{Source: arXiv:2608.26304v1 (CC BY 4.0)}
\begin{document}
\maketitle

\noindent\emph{Source and licence.} Multi-Objective Planning for Healthcare Facility Resilience: Mitigate Now or Respond Later?. Version arXiv:2608.26304v1 (CC BY 4.0), distributed by arXiv under the Creative Commons Attribution 4.0 International licence, http://creativecommons.org/licenses/by/4.0/, as recorded in the arXiv licence metadata for this exact version and shown on the abstract page https://arxiv.org/abs/2608.26304v1. Full licence notice: \texttt{licenses/arxiv-2608.26304-CC-BY-4.0.txt}.\par\smallskip

\noindent\textit{Editorial scope.} Counted unit: the article's proposition proposition (source0.tex lines 397--399). The article's complete original proof of it is delivered with it (source0.tex lines 401--455, one \texttt{proof} environment of 55 lines). No mathematics is written, repaired or re-derived: the statement and the proof are the article's own lines, in the article's own notation, and no retained line is substituted or re-flowed.  No further source-local context is needed: every cross-reference of every retained line resolves inside the two delivered spans.  Nothing was omitted from any retained span, so the package carries no omission record.  No retained line contains a citation, so the wrapper carries no bibliography and no bibliography entry.  No retained line contains a figure environment, an image-inclusion command or drawing code, so no image is delivered and the package declares no asset dependency.  Editorial formatting only, in addition: an AMS \texttt{amsart} preamble that restates the article's own declarations and macros used by the retained lines, namely the environments \texttt{proposition} on separate counters, together with the standard packages those lines need.  The retained environments print in this document as Proposition 1 in order of appearance, whereas the article prints the same items with the numbering of its own full document.  The complete native source of the article as distributed in the arXiv e-print for this exact version is bundled unchanged as \texttt{sources/208630/source0.tex}.
\begin{proposition}
The single-objective \(f_1\) version of the healthcare facility resilience planning problem is NP-hard.
\end{proposition}

\begin{proof}
Consider an arbitrary \(0\)--\(1\) knapsack instance with item weights \(w_j\),
profits \(\rho_j\), and capacity \(W\). Construct a restricted resilience
planning instance with a single scenario of probability one. For every sending
facility \(j\), let
\[
F_{sj}=H_j^{\max}=D_j=1,\qquad
C_j^H=w_j,
\]
and introduce one safe receiving facility \(k\) with sufficient capacity,
\(A_k\ge |J^{\mathrm{send}}|\). Set \(B=W\) and choose nonnegative restoration
and evacuation costs satisfying
\[
R_j+C_{jk}^E=\rho_j.
\]

Because \(F_{sj}=H_j^{\max}=1\), the hardening decision is binary:
\(y_j\in\{0,1\}\). If \(y_j=0\), the flood-status and evacuation constraints
force
\[
z_{sj}=\gamma_{sj}=q_{sjk}=1.
\]
If \(y_j=1\), setting
\[
z_{sj}=\gamma_{sj}=q_{sjk}=0
\]
is feasible and optimal because all associated costs are nonnegative.
Therefore, at optimality,
\[
z_{sj}=\gamma_{sj}=q_{sjk}=1-y_j.
\]
The restricted problem consequently becomes
\[
\min \sum_{j\in J^{\mathrm{send}}}
(R_j+C_{jk}^E)(1-y_j)
\qquad
\text{s.t.}
\qquad
\sum_{j\in J^{\mathrm{send}}} C_j^H y_j\le B,
\qquad
y_j\in\{0,1\}.
\]
Since \(\sum_j(R_j+C_{jk}^E)\) is constant, substituting the constructed
parameter values yields the equivalent problem
\[
\max \sum_{j\in J^{\mathrm{send}}}\rho_j y_j
\qquad
\text{s.t.}
\qquad
\sum_{j\in J^{\mathrm{send}}}w_j y_j\le W,
\qquad
y_j\in\{0,1\},
\]
which is the original \(0\)--\(1\) knapsack problem. This shows that the knapsack problem is a special case of the restricted instance of the resilience-planning problem. Since \(0\)--\(1\) knapsack is NP-hard, the restricted problem, and therefore the general facility resilience-planning problem, is also NP-hard.
\end{proof}

\end{document}
